%----------
%   IMPORTANTE
%----------

% Esta plantilla está basada en las recomendaciones de la guía "Trabajo fin de Grado: Escribir el TFG", que encontrarás en http://uc3m.libguides.com/TFG/escribir
% contiene recomendaciones de la Biblioteca basadas principalmente en estilos APA e IEEE, pero debes seguir siempre las orientaciones de tu Tutor de TFG y la normativa de TFG para tu titulación.



% ESTA PLANTILLA ESTÁ BASADA EN EL ESTILO APA


%----------
%	CONFIGURACIÓN DEL DOCUMENTO
%----------

\documentclass[12pt]{report} % fuente a 12pt

% MÁRGENES: 2,5 cm sup. e inf.; 3 cm izdo. y dcho.
\usepackage[
a4paper,
vmargin=2.5cm,
hmargin=3cm
]{geometry}

% INTERLINEADO: Estrecho (6 ptos./interlineado 1,15) o Moderado (6 ptos./interlineado 1,5)
\renewcommand{\baselinestretch}{1.15}
\parskip=6pt

% DEFINICIÓN DE COLORES para portada y listados de código
\usepackage[table]{xcolor}
\definecolor{azulUC3M}{RGB}{0,0,102}
\definecolor{gray97}{gray}{.97}
\definecolor{gray75}{gray}{.75}
\definecolor{gray45}{gray}{.45}

% Soporte para GENERAR PDF/A --es importante de cara a su inclusión en e-Archivo porque es el formato óptimo de preservación y a la generación de metadatos, tal y como se describe en http://uc3m.libguides.com/ld.php?content_id=31389625. 

% En la plantilla incluimos el archivo OUTPUT.XMPDATA. Puedes descargar este archivo e incluir los metadatos que se incorporarán al archivo PDF cuando compiles el archivo memoria.tex. Después vuelve a subirlo a tu proyecto.  
\usepackage[a-1b]{pdfx}

% ENLACES
\usepackage{hyperref}
\hypersetup{colorlinks=true,
	linkcolor=black, % enlaces a partes del documento (p.e. índice) en color negro
	urlcolor=blue} % enlaces a recursos fuera del documento en azul

% EXPRESIONES MATEMÁTICAS
\usepackage{amsmath,amssymb,amsfonts,amsthm}

% Codificación caracteres
\usepackage{txfonts} 
\usepackage[T1]{fontenc}
\usepackage[utf8]{inputenc}

% Definición idioma español
\usepackage[spanish, es-tabla]{babel} 
\usepackage[babel, spanish=spanish]{csquotes}
\AtBeginEnvironment{quote}{\small}

% diseño de PIE DE PÁGINA
\usepackage{fancyhdr}
\pagestyle{fancy}
\fancyhf{}
\renewcommand{\headrulewidth}{0pt}
\rfoot{\thepage}
\fancypagestyle{plain}{\pagestyle{fancy}}

% DISEÑO DE LOS TÍTULOS de las partes del trabajo (capítulos y epígrafes o subcapítulos)
\usepackage{titlesec}
\usepackage{titletoc}
\titleformat{\chapter}[block]
{\large\bfseries\filcenter}
{\thechapter.}
{5pt}
{\MakeUppercase}
{}
\titlespacing{\chapter}{0pt}{0pt}{*3}
\titlecontents{chapter}
[0pt]                                               
{}
{\contentsmargin{0pt}\thecontentslabel.\enspace\uppercase}
{\contentsmargin{0pt}\uppercase}                        
{\titlerule*[.7pc]{.}\contentspage}                 

\titleformat{\section}
{\bfseries}
{\thesection.}
{5pt}
{}
\titlecontents{section}
[5pt]                                               
{}
{\contentsmargin{0pt}\thecontentslabel.\enspace}
{\contentsmargin{0pt}}
{\titlerule*[.7pc]{.}\contentspage}

\titleformat{\subsection}
{\normalsize\bfseries}
{\thesubsection.}
{5pt}
{}
\titlecontents{subsection}
[10pt]                                               
{}
{\contentsmargin{0pt}                          
	\thecontentslabel.\enspace}
{\contentsmargin{0pt}}                        
{\titlerule*[.7pc]{.}\contentspage}  

% DISEÑO DE TABLAS y FIGURAS
\usepackage{multirow} % permite combinar celdas 
\usepackage{caption} % para personalizar el título de tablas y figuras
\usepackage{floatrow} % utilizamos este paquete y sus macros \ttabbox y \ffigbox para alinear los nombres de tablas y figuras de acuerdo con el estilo definido.
\usepackage{array} % con este paquete podemos definir en la siguiente línea un nuevo tipo de columna para tablas: ancho personalizado y contenido centrado
\newcolumntype{P}[1]{>{\centering\arraybackslash}p{#1}}
\DeclareCaptionFormat{upper}{#1#2\uppercase{#3}\par}
\usepackage{graphicx}
\graphicspath{{imagenes/}} % ruta a la carpeta de imágenes

% Diseño de tabla para ciencias sociales y humanidades
\captionsetup*[table]{
	justification=raggedright,
	labelsep=newline,
	labelfont=small,
	singlelinecheck=false,
	labelfont=bf,
	font=small,
	textfont=it
}

% Diseño de figuras para ciencias sociales y humanidades
\captionsetup[figure]{
	name=Figura,
	singlelinecheck=off,
	labelsep=newline,
	font=small,
	labelfont=bf,
	textfont=it
}
\floatsetup[figure]{
    style=plaintop,
    heightadjust=caption,
    footposition=bottom,
    font=small
}

% Configuración del pie de las figuras y tablas 
\captionsetup*[floatfoot]{
    footfont={small, up}
}

% NOTAS A PIE DE PÁGINA
\usepackage{chngcntr} % para numeración continua de las notas al pie
\counterwithout{footnote}{chapter}

% LISTADOS DE CÓDIGO
% soporte y estilo para listados de código. Más información en https://es.wikibooks.org/wiki/Manual_de_LaTeX/Listados_de_código/Listados_con_listings
\usepackage{listings}

% definimos un estilo de listings
\lstdefinestyle{estilo}{ frame=Ltb,
	framerule=0pt,
	aboveskip=0.5cm,
	framextopmargin=3pt,
	framexbottommargin=3pt,
	framexleftmargin=0.4cm,
	framesep=0pt,
	rulesep=.4pt,
	backgroundcolor=\color{gray97},
	rulesepcolor=\color{black},
	%
	basicstyle=\ttfamily\footnotesize,
	keywordstyle=\bfseries,
	stringstyle=\ttfamily,
	showstringspaces = false,
	commentstyle=\color{gray45},     
	%
	numbers=left,
	numbersep=15pt,
	numberstyle=\tiny,
	numberfirstline = false,
	breaklines=true,
	xleftmargin=\parindent
}

\captionsetup*[lstlisting]{font=small, labelsep=period}
% fijamos el estilo a utilizar 
\lstset{style=estilo}
\renewcommand{\lstlistingname}{\uppercase{Código}}


%BIBLIOGRAFÍA 

% CONFIGURACIÓN PARA LA BIBLIOGRAFÍA APA
\usepackage[style=apa, backend=biber, natbib=true, hyperref=true, uniquelist=false, sortcites]{biblatex}
\DeclareLanguageMapping{spanish}{spanish-apa}

% Para sustituir % por y o e
\makeatletter
\DefineBibliographyExtras{spanish}{%
  \setcounter{smartand}{1}%
  \let\lbx@finalnamedelim=\lbx@es@smartand
  \let\lbx@finallistdelim=\lbx@es@smartand
}

% Si te sigue apareciendo "y" en vez de "e" cuando un apellido empieza por "I", puedes forzar que aparezca la "e" añadiendo esto antes del nombre del autor en el archivo referencias.bib: {\forceE{Inicial del nombre}}. Por ejemplo, Luis {\forceE{I}}añez debería aparecer "e Iañez, L." si es el último autor.

\renewbibmacro*{name:delim:apa:family-given}[1]{%
  \ifnumgreater{\value{listcount}}{\value{liststart}}
    {\ifboolexpr{
       test {\ifnumless{\value{listcount}}{\value{liststop}}}
       or
       test \ifmorenames
     }
       {\printdelim{multinamedelim}}
       {\lbx@finalnamedelim{#1}}}
    {}}
\makeatother


% Añadimos las siguientes indicaciones para mejorar la adaptación del estilos en español
\DefineBibliographyStrings{spanish}{%
	andothers = {et\addabbrvspace al\adddot}
}
\DefineBibliographyStrings{spanish}{
	url = {\adddot\space[En línea]\adddot\space Disponible en:}
}
\DefineBibliographyStrings{spanish}{
	urlseen = {Acceso:}
}
\DefineBibliographyStrings{spanish}{
	pages = {pp\adddot},
	page = {p.\adddot}
}

\addbibresource{referencias.bib} % llama al archivo referencias.bib en el que deberá estar la bibliografía utilizada


%-------------
%	DOCUMENTO
%-------------

\begin{document}
\pagenumbering{roman} % Se utilizan cifras romanas en la numeración de las páginas previas al cuerpo del trabajo
	
%----------
%	PORTADA
%----------	
\begin{titlepage}
	\begin{sffamily}
	\color{azulUC3M}
	\begin{center}
		\begin{figure}[H] %incluimos el logotipo de la Universidad
			\makebox[\textwidth][c]{\includegraphics[width=16cm]{logo_UC3M.png}}
		\end{figure}
		\vspace{2.5cm}
		\begin{Large}
			Bachelor’s Degree Final Thesis\\			
			 2023-2024\\ %Indica el curso académico
			\vspace{2cm}		
			\textsl{Trabajo Fin de Grado}
			\bigskip
			
		\end{Large}
		 	{\Huge Modelling Opinion polarization in Social Networks through Markov Bridge Process}\\
		 	\vspace*{0.5cm}
	 		\rule{10.5cm}{0.1mm}\\
			\vspace*{0.9cm}
			{\LARGE Teresa Mattil Olazabal}\\ 
			\vspace*{1cm}
		\begin{Large}
			Tutor\\
			Iñaki Ucar\\
			UC3M Leganés, 12 de Julio 2024\\
		\end{Large}
	\end{center}
	\vfill
	\color{black}
	\end{sffamily}
\end{titlepage}

\newpage %página en blanco o de cortesía
\thispagestyle{empty}
\mbox{}

%----------
%	RESUMEN Y PALABRAS CLAVE
%----------	
\renewcommand\abstractname{\large\bfseries\filcenter\uppercase{Abstract}}
\begin{abstract}
\thispagestyle{plain}
\setcounter{page}{3}
	
	% ESCRIBIR EL RESUMEN AQUÍ
	
	\textbf{Keywords:}
	% Escribir las palabras clave aquí
        Eco chamber, polarization, Markov Bridge, HMM, Bayesian filtering
	
	\vfill
\end{abstract}
	\newpage % página en blanco o de cortesía
	\thispagestyle{empty}
	\mbox{}


%----------
%	DEDICATORIA
%----------	
\chapter*{Dedicatoria} % \chapter* evita que aparezca en el índice

\setcounter{page}{5}
	
	% ESCRIBIR LA DEDICATORIA AQUÍ	
		
	\vfill
	
	\newpage % página en blanco o de cortesía
	\thispagestyle{empty}
	\mbox{}
	

%----------
%	ÍNDICES
%----------	

%--
% Índice general
%-
\tableofcontents
\thispagestyle{fancy}

\newpage % página en blanco o de cortesía
\thispagestyle{empty}
\mbox{}

%--
% Índice de figuras. Si no se incluyen, comenta las líneas siguientes
%-
\listoffigures
\thispagestyle{fancy}

\newpage % página en blanco o de cortesía
\thispagestyle{empty}
\mbox{}

%--
% Índice de tablas. Si no se incluyen, comenta las líneas siguientes
%-
\listoftables
\thispagestyle{fancy}

\newpage % página en blanco o de cortesía
\thispagestyle{empty}
\mbox{}


%----------
%	MEMORIA
%----------	
\clearpage
\pagenumbering{arabic} % numeración con números arábigos para el resto de la memoria.	

\chapter{Introduction}

La polarización política es un fenómeno que preocupa cada vez más, al mismo tempo que la preocupación de la adicción a las redes sociales y como estas están diseñadas para atraparte. Las redes se han convertido en una fuente de información que muchas personas, especialmente jóvenes sustituyen por generadores de información oficiales como periódicos. Al basar nuestro conocimiento en una herramienta que además nos muestra el contenido que queremos ver para mantenernos conectados, surge el fenómeno de Eco Chambers: vamos formando una comunidad basada en homofilia  y accedemos limitadamente a información fuera de esa comunidad.
Este tema se ha tratado de muchas formas. The paradox of Information  [1] modela el sistema de intercambio de información de las redes sociales con ecuaciones parciales, mostrando que el sistema tiende a una estabilidad polarizada o uniforme según los parámetros. Otros han tratado de modelar las redes con Modelos Basados en Agentes [2], [3]. De esta forma, se pueden probar reglas de control para evitar llegar a polarización. También se ha planteado si realmente existe el fenómeno de echo chamber [4], desevidenciandolo con encuestas. En contra, se han creado modelos para justificar la coexistencia de ambos planteamientos: que existan pero no sean perceptibles [5].
R.Luo en su paper Echo Chamber and Seggregation in Social Networks muestra una perpectiva nueva: cómo estimar la polarización actual de una red. Modelando los diferentes niveles de polarización como cadenas de Markov, podemos estimar el estado de polarización actual de una red utilizando un Bayesian filter. Sin embargo, una cadena de Markov no capta el efecto de la temporalidad de una Echo Chamber. Por ello, propone utilizar un modelo markov Bridge para construir un proceso anticipatorio. 
Con este fin, el autor demuestra que un filtro Bayesiano de Markov Bridge es capaz de estimar mejor la evolución de la polarización que un filtro Markov. Pero para construir una cadena Markov Bridge, es necesario asumir el estado final.  Este Trabajo de Fin de Grado propone comprobar que una cadena Bridge no filtra mejor que una cadena de Markov cuando el estado final no es el estimado.

Para ello, el trabajo se divide en tres partes.
1) Construir un ABMS para capturar la dinámica de las redes sociales y poder sacar diferentes evoluciones de polarización según ciertos parámetros
2)  Diseño de modelos Hidden Markov Model y Hidden Markov Bridge, que sean capaces de capturar el estado de polarización de la red utilizando únicamente una muestra de esta, haciendolo aplicable en proyectos Data Driven. Y comparación de error entre HMB filter y HMM filter.
3) Comparación de HMM filter and HMB filter utilizando diferentes estados.

\chapter{Construction of an ABMS model to capture the dynamics of social networks}

To obtain results simulating a Social Network, we propose the given Agent-Based Model, based on the idea proposed by Muhammed Al Atiqi \cite{atiqi} that news providers are important stakeholders in people’s idea formations.

Many studies have proposed Agent-Based Modeling to capture Opinion Dynamics on SNS \cite{sns}. In general, they simulate the opinion on the network changing from a user posting a publication. The model proposed by Muhammed Atiqi allowed us to have a robust model sensitive to the factor of the news provider sentiments.

\textbf{Spoken summary of how interactions work + homophily type arguments ...}

The main factors of this model are that users publish being influenced by friends’ opinions, called by \cite{atiqi} social validation, and homophilic interactions.

\section{Proposed framework:}

The network is an undirected weighted graph with a set of $N$ user nodes $Nodes = \{ \text{user}_1, \text{user}_2, \ldots, \text{user}_N \}$ and $M$ weighted edges $E = \{ e_1, \ldots, e_M \}$. Each node $\text{user}_i$ has the following parameters:
\begin{itemize}
    \item An opinion $\theta_i$ ranging between $[-1, 1]$.
    \item A state $\text{state}_i \in \{\text{Uninformed}, \text{Informed}, \text{Accepted}, \text{Shared}\}$.
    \item Exposure $\epsilon_i \in [0,1]$.
    \item Hostility $\gamma_i$.
    \item $D_1$: Threshold for accepting a Neighbor’s post.
    \item $D_2$: Threshold for accepting a News Provider post.
    \item $\mu$: Constant factor that affects the way the opinion evolves after having accepted a post.
\end{itemize}

Independent of the nodes, there is a set of news that are published for each discrete time $t = \{0, \ldots, T\}$, $News = \{N_1, \ldots, N_{T \times p}\}$ with $p$ being the number of news published at each time $t$. Each news has an associated sentiment $s_j \in [-1, 1]$. The distance between news $n_j$ sentiment $s_j$ and user’s opinion $\theta_i$ is given by:
\begin{equation*}
    \text{dist}(\theta_i, s_j) = \theta_i - s_j
\end{equation*}
Exposure of $\text{user}_i$ is given by $\epsilon_i(t) = \epsilon(t-1) + \text{dist}(\theta_k, s_j)$. Hostility $\gamma_i = \theta_i \sum_k \theta_k$ for all neighbors $\text{user}_k$ of $\text{user}_i$.

\textbf{The simulation works as follows:}

\begin{enumerate}
    \item The network is initialized using a Barabasi-Albert graph \cite{barabasi} with $N$ nodes and $m$ XXX.
        \begin{itemize}
            \item Each node is initialized with a random opinion $\theta$, state $s = \text{Uninformed}$, exposure XXX, hostility XXX, and fixed thresholds and $\mu$.
        \end{itemize}
    \item For each time $t \in \{0, \ldots, N\}$, a set of $q$ news is published.
        \begin{itemize}
            \item Each time a news $n_j$ is published, all the nodes are exposed and decide to share the news if the absolute distance $|\text{dist}(\theta_i, s_j)| < d_2$.
            \item For all $\text{user}_k$: if there exists $(\text{user}_i, \text{user}_k) \in E$ and $\text{State}(\text{user}_i) = \text{Shared}$, $\text{State}(\text{user}_k) = \text{Informed}$ if exposure $\epsilon$.
            \item For all $\text{user}_k$: if $\text{State}(\text{user}_k) = \text{Informed}$ and the propagation is repeated $k$ times for news propagation, $\text{d}_1$, $\text{State}(\text{user}_k) = \text{Accepted}$ and $\theta_k = \theta_k + \mu \times \text{dist}(\theta_k, s_j)$.
            \item For all $\text{user}_k$: if $\text{State}(\text{user}_k) = \text{Accepted}$ and $\gamma_i > 0$ then $\text{State}(\text{user}_k) = \text{Shared}$.
        \end{itemize}
    \item The propagation is repeated $k$ times for news propagation. At the end of each news publication, the node’s states are reinitialized to uninformed.
    \item At the end of each time $t$, the state of the network is stored.
\end{enumerate}

\section{Polarization measure:}

We have considered some different polarization measures, including:

\begin{itemize}
    \item Algebraic connectivity of the graph $G$(t), denoted $\xi(G)$, is measure of how well-connected or resilient the graph is to disconnection. It is given by the second-smallest eigenvalue of the Laplacian matrix of $G$.
        \begin{equation*}
            \xi(G) = \lambda_2(L)
        \end{equation*}

    \item The graph conductance $\phi(G)$ of a graph $G$ is defined as the ratio of the number of edges leaving the smallest subset of vertices to the total number of edges in the graph. It is given by:
        \[
        \phi(G) = \frac{| \{ (u,v) \in E : u \in S, v \notin S \} |}{| E |}
        \]
        
        Where: $E$ is the set of edges in the graph, $S$ is the smallest subset of vertices in the graph and $| \cdot |$ denotes the cardinality of a set.
        
    \item Polarization score $\phi(t)$ is the ratio between interideological connections and intraideological connections of a node \cite{polarization}.
    \begin{equation*}
        \Phi(t) = \frac{\left|\{(\text{user}_i, \text{user}_j) | \theta_i \cdot \theta_j > 0\}\right|}{\left|\{(\text{user}_i, \text{user}_j) | \theta_i \cdot \theta_j < 0\}\right| + \left|\{(\text{user}_i, \text{user}_j) | \theta_i \cdot \theta_j > 0\}\right|}
    \end{equation*}
\end{itemize}

\section{Model Parameters}

We have studied the effects of different model parameters on the behavior of our system, including:

\begin{itemize}
    \item Different algorithms for graph initialization (complete, random, small world…)
    \item XXX
\end{itemize}

% ----------------------------------------------------------------------------------------------------------------------------------------------------------------------------------------------------------------------------------------------------------------------


\chapter{Evolution of polarization: HMM and HMB approach}
\section{Hidden Markov Bridge Model for Evolution of the Polarization of a SNS}

\subsection{Model Definition}

\subsubsection{Definition of Markov Process}
A Markov process is a stochastic process that satisfies the Markov property, which is to say, the probability of moving to the next state depends only on the present state and not on the sequence of events that preceded it.

\subsubsection{Definition of Markov Bridge Process}
A Markov Bridge process is a Markov process that has fixed starting and ending points, thus "bridging" the two end states over a specified time interval.

\subsubsection{Definition of Hidden Markov Model}
A Hidden Markov Model (HMM) is a statistical model in which the system being modeled is assumed to be a Markov process with unobserved (hidden) states.

Suppose that the MB process $\Phi = \{\phi(1), \ldots, \phi(t)\}$ is observed via the observation process $\tilde{\Phi} = \{\tilde{\phi}(1), \ldots, \tilde{\phi}(t)\}$. Assume that the observation at time $t$ given the state $\phi(t)$ is conditionally independent of $\phi(\tau)$ and $\tilde{\phi}(\tau)$, for $\tau \neq t$. This conditional independence implies that
\begin{equation}
P(\tilde{\phi}(1), \ldots, \tilde{\phi}(t)|\phi(1), \ldots, \phi(t)) = \prod_{k=1}^{t} P(\tilde{\phi}(k)|\phi(k)).
\end{equation}
The process $\tilde{\Phi}$ is called an HMB because the property (1) is analogous to the assumption made for HMM.

\subsection{Training the Model}

The transition probability of the first MB going from state $S[a]$ to state $S[b]$ is obtained by applying the Bayes rule as follows [25]:
\begin{equation}
B^{c}_{a,b}(t) = \frac{P_{a,b} \cdot P^{(T-(t+1))}_{b,c}}{P^{(T-t)}_{a,c}}
\end{equation}

\subsubsection{Proof of Bayes:}

\subsection{Baum-Welch Algorithm}
To obtain the transition matrix $P$, we use the Baum-Welch Smoothing algorithm for parameter tuning given a set of observations $\{\tilde{\phi}(1), \ldots, \tilde{\phi}(t)\}$.

\paragraph{Definition of Smoothing Algorithm}
A smoothing algorithm is an algorithm used to smooth out fluctuations in data, particularly in the context of signal processing and time series analysis.

\paragraph{Baum-Welch Algorithm Description}
The Baum-Welch algorithm is a particular case of a generalized expectation-maximization (EM) algorithm that can find the unknown parameters of an HMM.

\subsection{Bayes Filtering}

\subsubsection{Definition of HMM Bayes Filter}
An HMM Bayes filter is an algorithm used to compute the posterior distribution of the hidden states of an HMM given the observations.

The posterior probability can be evaluated recursively via Bayes’ rule as follows:
\begin{equation}
q_j(t+1) = O \ast \sum \left( \frac{B \ast q(t)}{O \ast B \ast q(t)} \right)
\end{equation}




\chapter{Setup and Results}


%----------
%	BIBLIOGRAFÍA
%----------	

%\nocite{*} % Si quieres que aparezcan en la bibliografía todos los documentos que la componen (también los que no estén citados en el texto) descomenta está línea

\clearpage
\addcontentsline{toc}{chapter}{Bibliografía}
\setquotestyle[english]{british} % Cambiamos el tipo de cita porque en el estilo IEEE se usan las comillas inglesas.
\printbibliography



%----------
%	ANEXOS
%----------	

% Si tu trabajo incluye anexos, puedes descomentar las siguientes líneas
%\chapter* {Anexo x}
%\pagenumbering{gobble} % Las páginas de los anexos no se numeran



\end{document}